\documentclass[12pt,a4paper]{article}
\usepackage[top=2.54cm,bottom=2.54cm,left=2.54cm,right=2.54cm]{geometry}
\usepackage{fontspec}
\usepackage{xeCJK}
\usepackage{array}
\usepackage{booktabs}
\usepackage{longtable}
\usepackage{graphicx}
\usepackage{hyperref}
\usepackage{amsmath}
\usepackage{listings}
\usepackage{setspace}
\usepackage{indentfirst}
\usepackage{xcolor}

\setmainfont{Times New Roman}
\setsansfont{Arial}
\setmonofont{Menlo}
\setCJKmainfont[BoldFont={Heiti TC}]{Songti TC}
\setCJKsansfont{Heiti TC}
\setCJKmonofont{Heiti TC}
\XeTeXlinebreaklocale "zh"
\XeTeXlinebreakskip = 0pt plus 1pt
\hypersetup{hidelinks}
\sloppy
\setlength{\parindent}{2em}
\setlength{\parskip}{0pt}
\onehalfspacing
\renewcommand{\contentsname}{目錄}

\newcommand{\Sinput}{S_{\mathrm{input}}}
\newcommand{\Svalidation}{S_{\mathrm{validation}}}
\newcommand{\Vtarget}{V_{\mathrm{target}}}
\newcommand{\Vpseudo}{V_{\mathrm{pseudo}}}
\newcommand{\Oroom}{\Omega_{\mathrm{room}}}
\newcommand{\Oocc}{\Omega_{\mathrm{occ}}}
\newcommand{\Ofree}{\Omega_{\mathrm{free}}}

\begin{document}
\pagenumbering{roman}

\begin{center}
{\LARGE\bfseries 國立彰化師範大學\\[0.6em]資訊工程學系碩士班\\[0.6em]碩士論文 OpenSpec 同步稿}
\end{center}

\begin{center}
{\Large\bfseries 單房間非連網家電環境影響學習之稀疏感測空間數位孿生原型}
\end{center}

\begin{center}
A Sparse-Sensing Spatial Digital Twin for Learning Environmental Impacts of Non-Networked Appliances in a Single Room
\end{center}

\begin{center}研究生：林昀佑\end{center}

\begin{center}指導教授：易昶霈 教授、沈慧宇 副教授\end{center}

\begin{center}版本：中文完整稿 v2.0，OpenSpec 同步版\end{center}

\begin{center}日期：2026 年 7 月 12 日\end{center}

\clearpage

\section*{審定書}
\addcontentsline{toc}{section}{審定書}

國立彰化師範大學資訊工程學系碩士班

碩士論文審定書

單房間非連網家電環境影響學習之稀疏感測空間數位孿生原型

研究生：林昀佑

本論文業經審查及口試合格，特此證明。

論文考試委員會召集人：

委員：

委員：

指導教授：易昶霈 博士

共同指導教授：沈慧宇 副教授

所長：

中華民國 115 年 月

\clearpage

\section*{誌謝}
\addcontentsline{toc}{section}{誌謝}

本研究能夠完成，首先感謝指導教授易昶霈教授與沈慧宇副教授在研究方向、方法與寫作上的指導與支持，以及各位口試委員的指正與建議。感謝求學過程中協助我的各位師長與同學，也感謝家人的支持與包容。

林昀佑 謹誌於

國立彰化師範大學資訊工程學系（所）

中華民國 115 年

\clearpage

\section*{摘要}
\addcontentsline{toc}{section}{摘要}

智慧建築與智慧居家系統需要掌握室內環境狀態，才能支援舒適度評估、能源管理與設備操作建議。然而，一般房間中的冷氣、電風扇、照明、窗戶與家具通常不會以完整建築管理系統形式提供高解析度資料；室內可部署的感測器數量也有限，難以直接取得真實完整三維連續場。本研究以單一臥室為場域，提出一個家具感知的稀疏感測空間數位孿生原型，針對溫度、相對濕度與照度建立可解釋、可校正且可驗證的目標點估計流程。本研究不以「少量感測器精確重建整個真實房間完整三維場」作為主要主張，而是將研究問題收斂為：在家具佔據區被排除後的自由空間中，如何利用模型允許使用的輸入感測器，估計枕頭、書桌、房間中央與家具邊界等目標位置的三因子狀態，並以保留感測器進行量化驗證。

方法上，本研究定義房間空間 $\Oroom$、家具佔據空間 $\Oocc$ 與可估計自由空間 $\Ofree=\Oroom\setminus\Oocc$，並將節點角色區分為 $\Sinput$、$\Svalidation$、$\Vtarget$ 與 $\Vpseudo$。其中 $\Sinput$ 可用於設備影響強度校正、低階空間 residual correction 與 residual learning；$\Svalidation$ 在模型擬合、正規化與模型選擇階段禁止使用，只在預測完成後作為 target-point validation truth。主模型採用變數專屬 reduced-order nominal model：溫度描述熱交換與局部熱源，濕度描述水氣交換與除濕近似，照度描述光源、窗戶日照、遮蔽與反射近似。此模型再與 IDW baseline、trilinear correction、hybrid residual correction 及後續 free-space modular estimators 進行比較。實體感測節點設計採 ESP32-C3、DHT11 與 BH1750 建立低成本三因子 sensing node；因房間含有床、書桌、櫃子與電風扇，建議部署由最小 6 顆擴充至 8--10 顆以上，並保留至少 2 顆 validation nodes。

評估採分層證據設計，分別區分 synthetic full-field、synthetic target holdout、real target-point、public task-aligned benchmark 與 intervention validation。既有受控模擬結果顯示，8 組標準情境中 base model 的平均 field MAE 為溫度 0.0474、濕度 0.1765、照度 2.0269，低於 IDW baseline 的 0.1723、0.4633、54.9052；hybrid residual leave-one-scenario-out 平均進一步降至 0.0017、0.0059、0.1407。既有 7 天 real-bedroom snapshot 中，pillow 參考點校正後 MAE 由 0.8967$^\circ$C、4.1286\% 與 309.0142 lux 降至 0.1676$^\circ$C、0.3939\% 與 16.6450 lux；此結果僅作為 pillow target-point evidence，不外推為完整真實三維場驗證。公開資料集 SML2010 與 CU-BEMS 僅作 task-aligned benchmark：SML2010 的 24 個任務中有 12 項取得最低 MAE，CU-BEMS 的 12 個任務中有 9 項優於 linear regression 但 0 項優於 persistence。本研究亦將人體存在、電風扇狀態與 Google Home UI 操作紀錄納入資料條件設計：occupancy 與 fan-on/off 時段須分組標記，Google Home UI 只作為 operator-verified operation event，不取代環境量測與 validation truth。推薦動作目前定位為 counterfactual action ranking，實際因果改善仍需後續 before/after 介入實驗驗證。

關鍵字：空間數位孿生、稀疏感測、家具感知、目標點驗證、非連網家電、溫度、濕度、照度、電風扇、Google Home。

\clearpage

\section*{Abstract}
\addcontentsline{toc}{section}{Abstract}

Smart building and smart home systems require indoor environmental awareness for comfort assessment, energy management, and device-operation recommendations. In ordinary rooms, however, air conditioners, fans, lights, windows, furniture, and occupant activities rarely provide complete building-management-system telemetry, and only a limited number of sensors can be deployed. This thesis proposes a furniture-aware sparse-sensing spatial digital twin prototype for a single bedroom. Rather than claiming that a small number of sensors can reconstruct the complete true three-dimensional indoor field, the thesis focuses on verifiable target-point estimation in the free air space after excluding furniture-occupied volumes.

The method defines the room domain $\Oroom$, the occupied furniture domain $\Oocc$, and the estimable free-space domain $\Ofree=\Oroom\setminus\Oocc$. Physical nodes are separated into $\Sinput$ and $\Svalidation$, while research query locations and model-generated support values are represented as $\Vtarget$ and $\Vpseudo$. Only $\Sinput$ may be used for calibration, correction, and residual learning; $\Svalidation$ is held out until prediction is complete and is used only for target-point validation. The model uses variable-specific reduced-order nominal structures for temperature, relative humidity, and illuminance, and compares the interpretable physics-based estimator with IDW, trilinear residual correction, hybrid residual learning, and future modular free-space estimators. The planned hardware layer uses ESP32-C3, DHT11, and BH1750 to build low-cost three-factor sensing nodes. Because the bedroom includes furniture and an electric fan, the deployment plan is expanded beyond a pure 8-corner assumption toward an 8--10 node furniture-aware layout with at least two validation nodes.

The evidence design separates synthetic full-field, synthetic target-holdout, real target-point, public task-aligned benchmark, and intervention-validation evidence. Existing controlled simulation results show that the base model obtains average field MAE of 0.0474, 0.1765, and 2.0269 for temperature, humidity, and illuminance, compared with 0.1723, 0.4633, and 54.9052 for IDW. Leave-one-scenario-out hybrid residual correction further reduces average MAE to 0.0017, 0.0059, and 0.1407. A seven-day real-bedroom pillow snapshot reduces pillow-point error from 0.8967$^\circ$C, 4.1286\%, and 309.0142 lux to 0.1676$^\circ$C, 0.3939\%, and 16.6450 lux, but this is interpreted only as pillow target-point evidence rather than dense 3-D ground truth. Public SML2010 and CU-BEMS results are used only as task-aligned external benchmarks. Occupancy, fan state, and Google Home UI operation logs are treated as experimental context: fan-on/off and occupied/unoccupied periods must be separated, and Google Home UI events are treated as operator-verified operation events rather than environmental truth. Action recommendations remain counterfactual rankings and require future before/after intervention studies for causal claims.

Keywords: spatial digital twin, sparse sensing, furniture-aware estimation, target-point validation, non-networked appliances, temperature, humidity, illuminance, electric fan, Google Home.

\clearpage

\tableofcontents
\clearpage
\pagenumbering{arabic}

\section{緒論}

\subsection{研究背景}

智慧家庭與智慧建築系統逐漸被用於室內環境監控、舒適度評估、能源管理與設備操作建議。這類系統若要判斷使用者所在位置是否過熱、過濕或過暗，不能只知道單一牆面溫度或單一燈光狀態，而需要理解空間中不同位置的溫度、相對濕度與照度差異。然而，一般臥室或小型房間很少具備完整建築管理系統，也不可能在每個空間點布建感測器；使用者真正關心的位置，例如枕頭、書桌、床邊或家具邊界，往往不是既有感測器所在位置。

另一方面，既有家電與環境作用源的資料可得性並不一致。冷氣可能可由 Google Home UI 操作，但未必能提供完整出風量、擺葉角度或實際冷卻效果；電風扇會改變氣流混合，卻不是冷氣式降溫設備；照明可能可被智慧開關控制，但照度仍受燈具方向、人體遮擋與家具遮蔽影響；窗戶、床、書桌與櫃子也會改變自由空間中的熱、濕與光分布。若數位孿生模型只依賴裝置 API 或只假設空房間幾何，便無法形成可防守的真實房間驗證。

\subsection{研究動機}

本研究原型初期曾以 8 顆角落感測器與簡化三因子場模型作為基線。此基線有助於建立可解釋的 physics-first 估計流程，但也暴露出三個限制。第一，固定 8-corner layout 只適合空矩形房間，無法反映床、櫃子、書桌等家具佔據空間。第二，若沒有保留 validation sensors，就容易把校正後的模型輸出誤當成真實未量測位置的 ground truth。第三，電風扇、occupancy 與 Google Home 操作紀錄若未被標記，會讓不同邊界條件下的資料混在一起，造成模型誤差解釋不清。

因此，本研究將核心定位從「8 顆角落感測器完整重建真實三維場」收斂為「家具感知自由空間中的可驗證目標點估計」。8-corner 模型保留為 baseline 與可證明起點；正式研究框架則加入 $\Sinput$ / $\Svalidation$ 角色分離、target-point holdout、家具自由空間、電風扇狀態分組、occupancy 標記，以及 Google Home UI 操作紀錄作為 operation context。

\subsection{研究問題}

本研究問題重新整理如下：

\begin{itemize}
\item RQ1：在相同 $\Sinput$ 與 $\Svalidation$ 切分下，加入變數專屬物理先驗的三因子估計方法，是否比純距離插值 IDW 更能降低目標點誤差？
\item RQ2：在房間含有家具遮蔽時，定義 $\Ofree=\Oroom\setminus\Oocc$ 並使用 furniture-aware sensing layout，是否比原始固定 8-corner 假設更合理地估計 pillow、desk、room center 與 near-furniture boundary 等目標位置？
\item RQ3：在不使用 $\Svalidation$ 的條件下，active-device power calibration、trilinear correction 與 hybrid residual learning 是否能降低保留目標點或未見情境的誤差？
\item RQ4：冷氣、電風扇、照明、occupancy 與 Google Home UI operation events 應如何被標記為 experimental context，使真實房間驗證不混淆不同邊界條件？
\item RQ5：在完成真實 before/after intervention 前，模型輸出的控制建議應如何限制為 counterfactual action ranking，而不宣稱已證明因果控制效果？
\end{itemize}

\subsection{研究範圍與限制}

本研究範圍與限制如下：

\begin{itemize}
\item 場域限定為單一臥室，不處理多房間或跨樓層空氣交換。
\item 研究核心為目標點三因子估計，不主張由稀疏節點取得完整真實三維連續場。
\item 家具佔據空間不作為一般空氣查詢點，主要估計區域為 $\Ofree$。
\item 實體感測節點為低成本研究 node；DHT11 不被定位為高精度實驗室 reference。
\item Google Home UI 操作紀錄可作為 operator-verified operation event，但不能取代溫度、濕度與照度量測。
\item 電風扇被視為 airflow redistribution source，而不是等同冷氣的主動降溫設備。
\item 人體存在被視為 dynamic heat/moisture/obstruction source；occupied 與 unoccupied periods 不應混為同一 validation condition。
\item MCP、Web demo 與 agent tools 為應用介面層，不作為主要科學貢獻。
\item 控制建議目前僅作候選動作排序，實際因果效果需另以 before/after intervention 驗證。
\end{itemize}

\subsection{預期貢獻}

本研究預期貢獻如下：

\begin{enumerate}
\item 建立單房間三因子稀疏感測數位孿生之 OpenSpec-driven 研究契約，明確區分實作、驗證、提案延伸與未來工作。
\item 提出 $\Sinput$ / $\Svalidation$ / $\Vtarget$ / $\Vpseudo$ 的資料角色設計，降低 target-point validation leakage 風險。
\item 將原始 8-corner baseline 擴充為 furniture-aware free-space estimation 問題，支援家具邊界、枕頭、書桌與房間中心等真實目標點。
\item 設計低成本 ESP32-C3 + DHT11 + BH1750 sensing node 與 8--10 顆以上的家具感知部署策略。
\item 將電風扇、occupancy 與 Google Home UI 操作紀錄納入 experimental context，使真實資料分析能分組處理不同邊界條件。
\item 以分層證據框架整合 synthetic full-field、synthetic target holdout、real target-point、public task-aligned benchmark 與 intervention validation。
\end{enumerate}

\section{文獻探討與研究定位}

\subsection{室內環境建模}

室內環境建模可大致分為高解析數值模擬、reduced-order / grey-box model、資料驅動模型與混合模型。CFD 可細緻描述流場、熱交換與局部對流，但需要完整幾何、材料與邊界條件，且運算成本高。相較之下，控制導向 reduced-order model 更適合用於即時估計與原型系統，雖犧牲部分物理細節，卻能保留可解釋性與校正能力。本研究採取此路線，並明確限制其證據範圍：模型產生的 dense grid 是 estimate，只有受控模擬 reference field 才能稱為 synthetic truth。

\subsection{空間插值與目標點估計}

IDW 等空間插值法是稀疏感測場估計的常見 baseline。其優點是不需設備先驗，缺點是無法表達冷氣出風、窗戶日照、燈具遮蔽、家具佔據與風扇混合等語意。因此，本研究將 IDW 作為 baseline，而不將其視為完整解法。後續的 2-D triangulation、3-D tetrahedral interpolation 與 Cell-IDW Fusion 被定位為 modular free-space estimators；在尚未完成實作與真實驗證前，必須標示為 proposed extension 或 future work。

\subsection{數位孿生與智慧居家操作紀錄}

數位孿生在智慧建築中通常整合感測資料、模型狀態與操作決策。本研究的特點在於場域並非完整 BMS，而是低成本單房間環境。Google Home UI 可提供冷氣、電風扇與照明的操作事件時間線，但即使操作成功，它仍只代表 device operation context，而不是環境 measurement。環境 truth 必須來自感測器，尤其是未參與模型擬合的 $\Svalidation$。

\subsection{occupancy 與 fan effect}

人在房間內會產生顯熱、潛熱與局部遮蔽，電風扇則會改變氣流混合與局部對流。這兩者都會造成感測資料分布變化。若將 occupied、unoccupied、fan-on 與 fan-off 直接混算，模型誤差將難以解釋。因此本研究將 occupancy 與 fan state 視為 experimental context，至少在資料層標記並分組分析。

\section{系統架構與數學模型}

\subsection{整體架構}

本研究架構分為六層：空間定義層、感測節點層、物理估計層、校正與 residual learning 層、驗證與證據層，以及應用介面層。空間定義層描述 $\Oroom$、$\Oocc$ 與 $\Ofree$；感測節點層描述 $\Sinput$ 與 $\Svalidation$；物理估計層對溫度、濕度與照度建立變數專屬 nominal models；校正層處理 active-device power calibration、trilinear correction 與 hybrid residual；驗證層保證 validation observations 不進入 fitting；應用層則提供 Web、MCP 或 agent tools。

\subsection{空間與節點角色}

房間空間定義為：

\begin{equation}
\Ofree = \Oroom \setminus \Oocc
\end{equation}

其中 $\Oroom$ 為房間總空間，$\Oocc$ 為家具佔據體積，$\Ofree$ 為可估計自由空氣空間。若 query point 位於家具內部，系統應回傳 occupied 或低可信度狀態，而不將其納入自由空間 error summary。

節點角色定義如下：

\begin{longtable}{p{0.22\linewidth}p{0.68\linewidth}}
\toprule
角色 & 定義 \\
\midrule
$\Sinput$ & 可被模型用於 calibration、correction、feature construction 與 residual fitting 的真實感測器。 \\
$\Svalidation$ & 模型擬合階段禁止使用，只在 prediction 完成後由 evaluator 讀取的保留感測器。 \\
$\Vtarget$ & 枕頭、書桌、房間中央、家具邊界等研究目標位置；不一定有實體感測器。 \\
$\Vpseudo$ & 模型或幾何方法產生的支撐點或補值，不是 ground truth。 \\
\bottomrule
\end{longtable}

集合關係為：

\begin{equation}
S_{\mathrm{all}} = \Sinput \cup \Svalidation, \qquad \Sinput \cap \Svalidation = \varnothing
\end{equation}

\subsection{三因子 nominal model}

本研究不使用同一公式描述三種環境量，而採變數專屬模型。

溫度估計可概念化為：

\begin{equation}
T(q,t) = T_{\mathrm{base}}(t) + \Delta T_{\mathrm{AC}}(q,t) + \Delta T_{\mathrm{window}}(q,t) + \Delta T_{\mathrm{light}}(q,t) + \Delta T_{\mathrm{occupant}}(q,t) + r_T(q,t)
\end{equation}

濕度估計可概念化為：

\begin{equation}
RH(q,t) = RH_{\mathrm{base}}(t) + \Delta RH_{\mathrm{AC}}(q,t) + \Delta RH_{\mathrm{window}}(q,t) + \Delta RH_{\mathrm{occupant}}(q,t) + r_{RH}(q,t)
\end{equation}

照度估計可概念化為：

\begin{equation}
L(q,t) = L_{\mathrm{base}}(t) + \Delta L_{\mathrm{light}}(q,t) + \Delta L_{\mathrm{window}}(q,t) - Shade_{\mathrm{furniture}}(q,t) - Shade_{\mathrm{occupant}}(q,t) + r_L(q,t)
\end{equation}

其中 $q$ 為查詢位置，$t$ 為時間，$r_T$、$r_{RH}$ 與 $r_L$ 為 residual terms。

\subsection{電風扇、occupancy 與操作事件}

電風扇不被建模為冷氣式降溫來源，而是 airflow redistribution source。其第一版處理方式是將 fan state 作為 experimental condition，至少記錄 power、speed、oscillation、direction 與是否位於主要風路徑。fan-on 與 fan-off 不混算；unknown fan state 不作主要 validation evidence，除非明確標示限制。

occupancy 被視為動態熱濕與遮蔽來源。正式 validation 應至少區分 unoccupied、occupied sleeping、occupied desk、occupied moving 與 unknown。若 occupancy unknown，該時段不應作為主要 ground-truth validation。

Google Home UI 操作紀錄被定位為 operation context。若所有冷氣、風扇與照明操作皆透過 Google Home UI 執行，且操作者確認操作成功，可將事件標記為 \texttt{ui\_operator\_verified} 或 \texttt{verified\_state}。但該紀錄仍不取代環境 sensing node 或 $\Svalidation$。

\subsection{感測器校正與 residual correction}

正式資料收集前，所有 sensing nodes 應進行 12--24 小時同位置校正，計算 temperature offset、humidity offset 與 light scale 或 offset。模型校正流程先使用 $\Sinput$ 估計 active-device power scale，再建立 trilinear residual correction 或 residual learning dataset。$\Svalidation$ 只能在 prediction 完成後進入 evaluator。

\section{實體感測節點與資料收集設計}

\subsection{低成本三因子 sensing node}

本研究 v1 實體節點採低成本組成：ESP32-C3 作為 Wi-Fi MCU，DHT11 作為低成本溫濕度感測器，BH1750 作為 digital lux sensor。節點不定位為高精度實驗室儀器，而是可大量部署的 sparse sensing node。DHT11 的限制需在論文中明確揭露；若 validation nodes 預算允許，可升級為 DHT22 或 SHT31。

\begin{longtable}{p{0.22\linewidth}p{0.28\linewidth}p{0.4\linewidth}}
\toprule
模組 & 建議元件 & 研究用途 \\
\midrule
MCU & ESP32-C3 & Wi-Fi、timestamp、payload upload。 \\
溫濕度 & DHT11 & 低成本 temperature / humidity input，需同位置校正。 \\
照度 & BH1750 & digital lux sensor，優先於未校正 LDR。 \\
固定方式 & 開放式基座或打孔外殼 & 固定位置、方向與 USB 線，避免密閉自熱。 \\
供電 & USB 5V & 正式長期資料收集先不採電池版。 \\
\bottomrule
\end{longtable}

\subsection{家具與風扇感知部署}

因房間有床、書桌、櫃子與電風扇，節點部署不應只依賴 8 個角落。建議最小 6 顆、可防守 8--10 顆、dense 12--14 顆。10 顆建議配置如下：

\begin{longtable}{p{0.34\linewidth}p{0.18\linewidth}p{0.38\linewidth}}
\toprule
節點 & 角色 & 目的 \\
\midrule
\texttt{input\_center} & input & 房間自由空間基準。 \\
\texttt{input\_window\_side} & input & 外氣、日照與窗邊梯度。 \\
\texttt{input\_ac\_path} & input & 冷氣氣流與降溫/除濕作用。 \\
\texttt{input\_fan\_path} & input & 電風扇主風路徑。 \\
\texttt{input\_fan\_shadow\_zone} & input & 家具後方或風較難到達區域。 \\
\texttt{input\_near\_furniture\_left} & input & 家具邊界一側。 \\
\texttt{input\_near\_furniture\_right} & input & 家具邊界另一側。 \\
\texttt{input\_bed\_side} & input & 床邊自由空間與人體影響比較。 \\
\texttt{validation\_pillow} & validation & 枕頭目標點 holdout truth。 \\
\texttt{validation\_desk} & validation & 書桌目標點 holdout truth。 \\
\bottomrule
\end{longtable}

\subsection{資料 payload 與 operation event}

sensing node payload 至少包含 node id、role、position、timestamp、temperature、humidity、illuminance、raw values 與 quality flags。operation event payload 另記錄 Google Home UI 或人工紀錄產生的設備操作事件。兩者不可混淆：sensing payload 是環境量測，operation event 是 device context。

operation event 建議欄位包含：

\begin{lstlisting}[basicstyle=\ttfamily\small,breaklines=true]
{
  "schema": "operation_event_v1",
  "room_id": "bedroom_01",
  "timestamp": "2026-07-12T21:30:01+08:00",
  "source": "google_home",
  "source_detail": "google_home_ui",
  "device_id": "fan_main",
  "device_type": "fan",
  "requested_state": {"power": "on", "speed": "medium"},
  "state_confidence": "ui_operator_verified",
  "verification_method": "ui_observed_success"
}
\end{lstlisting}

\section{系統實作與服務介面}

\subsection{Python 原型}

Python 原型負責房間、裝置、感測器、校正與估計器邏輯。核心實作應維持 role-aware：任何 calibration、trilinear fitting、residual dataset construction 或 model selection 都不得讀取 $\Svalidation$。這一點不僅是程式約束，也是論文證據有效性的前提。

\subsection{Estimator interface}

為避免每種方法使用不同資料切分，本研究規劃統一 estimator contract。每個 estimator 應接收相同 context、query point 與 $\Sinput$ observations，並回傳 value、method、support nodes、confidence 與 provenance。此設計可比較 BasePhysics、IDW、BasePhysics+trilinear、hybrid residual、2-D triangulation、3-D tetrahedral 與 Cell-IDW Fusion。

\subsection{MCP、Web demo 與 agent tools}

MCP tools、Web demo 與 agent tools 是系統應用與展示介面，讓使用者或 AI client 查詢環境狀態、執行反事實模擬或檢視推薦動作。這些介面不構成主要科學貢獻，也不能替代測量與驗證。論文中應將其歸類為 application layer。

\section{實驗設計與結果分析}

\subsection{證據分層}

本研究採取下列表格所示的證據分層，以避免不同證據類型支撐超出其能力範圍的主張。

\begin{longtable}{p{0.25\linewidth}p{0.34\linewidth}p{0.31\linewidth}}
\toprule
證據層 & 可支持主張 & 不可支持主張 \\
\midrule
Synthetic full-field & 在受控模擬 truth 下比較方法表現。 & 不能代表真實房間完整 3-D ground truth。 \\
Synthetic target holdout & 檢查 $\Sinput$ / $\Svalidation$ 分離與 target error。 & 不能取代真實感測器驗證。 \\
Real target-point & 支持特定 validation node 的實測目標點誤差。 & 不能外推到所有未量測位置。 \\
Public task-aligned benchmark & 檢查時間預測或 task-mapping 能力。 & 不能證明房間三維空間場準確。 \\
Intervention validation & 可檢驗 before/after 操作因果改善。 & 未執行前不能宣稱控制有效。 \\
\bottomrule
\end{longtable}

\subsection{既有 synthetic full-field 結果}

既有 8 組標準情境中，BasePhysics 在 synthetic full-field reference 下優於 IDW baseline。此結果支持「變數專屬物理先驗在受控模擬中可降低場估計誤差」，但不應被描述為真實房間完整 3-D field 驗證。

\begin{longtable}{p{0.28\linewidth}ccc}
\toprule
方法 & 溫度 MAE & 濕度 MAE & 照度 MAE \\
\midrule
IDW baseline & 0.1723 & 0.4633 & 54.9052 \\
BasePhysics & 0.0474 & 0.1765 & 2.0269 \\
LOO Hybrid residual & 0.0017 & 0.0059 & 0.1407 \\
\bottomrule
\end{longtable}

\subsection{既有 real-bedroom pillow snapshot}

既有 7 天 real-bedroom snapshot 顯示，pillow 參考點校正後誤差明顯下降。新版論文應將此結果標記為 pillow target-point evidence，而非完整房間場驗證。

\begin{longtable}{p{0.28\linewidth}ccc}
\toprule
狀態 & 溫度 MAE & 濕度 MAE & 照度 MAE \\
\midrule
校正前 & 0.8967$^\circ$C & 4.1286\% & 309.0142 lux \\
校正後 & 0.1676$^\circ$C & 0.3939\% & 16.6450 lux \\
\bottomrule
\end{longtable}

\subsection{public task-aligned benchmark}

公開資料集 SML2010 與 CU-BEMS 不是本研究臥室三維空間真值。其用途是檢查模型特徵映射與時間任務相容性。既有結果顯示，SML2010 的 24 個任務中有 12 項取得最低 MAE；CU-BEMS 的 12 個任務中有 9 項優於 linear regression，但沒有優於 persistence。論文必須保留 persistence 優於本模型的結果，並說明其原因可能來自 CU-BEMS 任務的強時間自相關與資料任務性質，而不是隱藏不利結果。

\subsection{fan、occupancy 與 Google Home 分組分析}

真實房間資料必須額外分組：fan-on / fan-off、occupied / unoccupied、AC on/off、light on/off、window open/closed。Google Home UI 操作成功紀錄可用來建立 operation event timeline，但各組環境結果仍以 sensing node readings 計算。fan unknown 或 occupancy unknown 的時段不應進入主要 validation summary，除非明確標示 uncertainty。

\section{控制推薦、限制與未來工作}

\subsection{counterfactual action ranking}

本研究可根據模型預測輸出候選操作排序，例如冷氣開啟、風扇開啟、照明調整或窗戶開啟程度。但在尚未進行真實 before/after intervention 前，這些輸出只能稱為 counterfactual action ranking。其格式應包含 current state、candidate action、predicted delta、comfort penalty、uncertainty 與 provenance。

\subsection{研究限制}

目前限制如下：

\begin{itemize}
\item 真實驗證仍以少數 target points 為主，不能代表全室所有點。
\item DHT11 精度有限，validation nodes 若作正式論文結果，需標明校正流程與感測器限制。
\item 電風扇與 occupancy 尚未完成完整物理建模，第一版主要透過 context labeling 與分組分析處理。
\item Google Home UI operation event 高可信度代表操作成功，但不代表環境效果已達穩態。
\item 2-D triangulation、3-D tetrahedral 與 Cell-IDW Fusion 若尚未完成實作與實測，需標記為 proposed extension。
\item intervention validation 尚未完成前，不得宣稱推薦動作具有已證明的因果改善效果。
\end{itemize}

\subsection{未來工作}

未來工作包括：

\begin{enumerate}
\item 完成 bedroom\_01 的實際 node deployment map，記錄 node id、role、$x/y/z$、height、light orientation、nearby furniture 與 fan-relative position。
\item 建立 firmware schema、MQTT topic validation tests 與 operation event validation tests。
\item 執行 12--24 小時同位置校正並產生 offset / scale 計算腳本。
\item 蒐集多日 real target-point data，至少涵蓋 pillow、desk、room center 與 near-furniture boundary。
\item 實作統一 estimator interface 與 Cell-IDW / triangulation / tetrahedral estimators。
\item 建立 claim-to-evidence matrix，將中文論文、IEEE 稿、PPT、speaker notes、figures 與 output JSON 同步。
\item 設計 before/after intervention protocol，以檢驗 counterfactual action ranking 的實際因果效果。
\end{enumerate}

\section{結論}

本研究提出一個以單房間為場域的家具感知稀疏感測三因子空間數位孿生原型。相較於將少量角落感測器描述為完整三維場重建，本研究將主張收斂為自由空間中的可驗證目標點估計，並以 $\Sinput$ / $\Svalidation$ 分離降低驗證洩漏風險。既有 synthetic 結果支持變數專屬物理先驗與 residual correction 在受控場中的效果，既有 pillow snapshot 提供特定 real target-point evidence，公開資料集則只作 task-aligned benchmark。實體部署層面，本研究規劃使用 ESP32-C3、DHT11 與 BH1750 建立低成本 sensing nodes，並因家具與電風扇存在而採用 8--10 顆以上的家具感知部署。occupancy、fan state 與 Google Home UI operation events 被納入 experimental context，使未來真實資料分析可以分組、追溯且不混淆證據邊界。

\clearpage

\section*{參考文獻}
\addcontentsline{toc}{section}{參考文獻}

[1] ASHRAE, \textit{ASHRAE Handbook: Fundamentals}. Atlanta, GA: ASHRAE.

[2] T. A. Reddy, \textit{Applied Data Analysis and Modeling for Energy Engineers and Scientists}. Springer.

[3] S. Wang and X. Xu, ``Simplified building model for transient thermal performance estimation using GA-based parameter identification,'' \textit{International Journal of Thermal Sciences}.

[4] F. Jazizadeh, G. Kavulya, L. Klein, and B. Becerik-Gerber, ``Continuous sensing of occupant perception of indoor ambient factors,'' \textit{Automation in Construction}.

[5] H. Oh et al., ``Hybrid modeling based on integrating simulation and operational data to improve indoor air temperature predictions, a controlled variable in digital twin models,'' 2024.

[6] OpenSpec project files in this repository: research-contract, spatial-estimation, evidence-and-artifacts, and sensing-node specifications.

\clearpage

\appendix

\section{OpenSpec 同步檢查摘要}

本稿已同步下列 OpenSpec 研究邊界：

\begin{itemize}
\item 核心主張改為家具感知自由空間中的 target-point estimation。
\item 明確加入 $\Sinput$、$\Svalidation$、$\Vtarget$ 與 $\Vpseudo$。
\item 保留 8-corner 作為 baseline，而不是唯一正式部署假設。
\item 加入 sensing node：ESP32-C3、DHT11、BH1750。
\item 加入 8--10 顆以上家具與風扇感知部署策略。
\item 加入 fan-on/off、occupancy 與 Google Home UI operation context。
\item 將 public datasets 限定為 task-aligned benchmark。
\item 將 recommendation 限定為 counterfactual action ranking。
\end{itemize}

\section{待補產物同步}

後續仍需同步：

\begin{itemize}
\item IEEE 英文稿。
\item PPT 與 speaker notes。
\item claim-to-evidence matrix。
\item 實際 bedroom\_01 node deployment map。
\item 真實多日 target-point validation figures。
\item output JSON 與 thesis figures 的數值一致性檢查。
\end{itemize}

\end{document}
